% Part: computability
% Chapter: recursive-functions
% Section: non-pr-functions

\documentclass[../../../include/open-logic-section]{subfiles}

\begin{document}

\olfileid[fr]{cmp}{rec}{npr}
\olsection{Fonctions non récursives primitives}

Les fonctions récursives primitives n'épuisent pas les fonctions
calculables au sens intuitif. Il paraît clair que nous pouvons
énumérer toutes les fonctions récursives primitives unaires sous
la forme $f_0,f_1,f_2,~\dots$, de telle sorte que la valeur de
$f_x$ en~$y$ soit calculable effectivement. Autrement dit, la
fonction $g(x,y)$ définie par
\[
g(x,y) = f_x(y)
\]
est calculable. Il en va donc de même de la fonction
\begin{eqnarray*}
h(x) & = & g(x,x) + 1 \\
& = & f_x(x) +1.
\end{eqnarray*}
Pour chaque fonction récursive primitive $f_i$, les valeurs de
$h$ et de $f_i$ diffèrent en~$i$. Ainsi, $h$ est calculable
mais n'est pas récursive primitive; il en va de même de~$g$.
C'est une version « effective » de l'argument diagonal de Cantor.

On peut donner des exemples plus explicites de fonctions calculables
qui ne sont pas récursives primitives. Ainsi, notons $g^n(x)$
l'expression $g(g(\dots g(x)))$, où $g$ apparaît $n$~fois, avec $g^0(x)=x$, et
définissons une suite de fonctions $g_0,g_1,\dots$ par
\begin{eqnarray*}
g_0(x) & = & x+1 \\
g_{n + 1}(x) & = & g_n^x(x)
\end{eqnarray*}
On peut vérifier que chaque fonction $g_n$ est récursive primitive.
Chacune croît beaucoup plus vite que la précédente :
$g_1(x)=2x$, $g_2(x)=2^x\cdot x$, et la croissance de
$g_3(x)$ est approximativement celle d'une tour exponentielle
formée de $x$~copies de~$2$. La fonction d'Ackermann--Péter
correspond essentiellement à la fonction $G(x)=g_x(x)$;
on peut montrer que celle-ci finit par dominer toute fonction
récursive primitive.

Revenons à l'énumération des fonctions récursives primitives.
Rappelons que nous avons attribué une notation symbolique à chacune
d'elles : il suffit donc d'énumérer ces notations. À chaque
notation~$F$, nous pouvons associer récursivement un entier
naturel $\#(F)$ de la manière suivante :
\begin{eqnarray*}
\#(0) & = & \langle 0 \rangle \\
\#(S) & = & \langle 1 \rangle \\
\#(\Proj{n}{i}) & = & \langle 2, n, i \rangle \\
\#(\fn{Comp}_{k,l}[H,G_0,\dots,G_{k-1}]) & = & \langle
3,k,l,\#(H),\#(G_0),\dots,\#(G_{k-1}) \rangle \\
\#(\fn{Rec}_l[G,H]) & = & \langle 4, l, \#(G), \#(H) \rangle
\end{eqnarray*}
Nous utilisons ici le codage des suites de nombres par des entiers
naturels présenté dans la section précédente. Chaque notation
reçoit ainsi un code numérique. Certaines suites, et donc certains
entiers, ne correspondent certes à aucune notation. Posons alors
$f_i$ égal à la fonction récursive primitive unaire dont la
notation porte le code~$i$, si $i$ code une telle notation, et
égal à la fonction constante nulle dans le cas contraire. Nous
obtenons ainsi une énumération explicite de toutes les fonctions
récursives primitives unaires.

(Certaines fonctions, dont la fonction constante nulle, figurent
plusieurs fois dans cette liste. Cela ne tient pas seulement à
notre codage : la fonction constante nulle possède elle-même
plusieurs notations. Il existe pourtant des énumérations effectives
sans répétition. En revanche, aucune fonction calculable
ne décide si une notation donnée représente
la fonction constante nulle.)\footnote{Correction éditoriale : l'impossibilité générale d'une énumération effective sans répétition, affirmée dans la source, est fausse. Voir Shih-Chao Liu, \emph{An enumeration of the primitive recursive functions without repetition}, \emph{Tohoku Mathematical Journal}, deuxième série, 12(3), 1960, p.~400–402, \href{https://doi.org/10.2748/tmj/1178244403}{doi:10.2748/tmj/1178244403}. Ce résultat ne fournit pas de décision d'égalité à la fonction nulle à partir d'une notation arbitraire.}

Nous pouvons maintenant poser $g(x,y)=f_x(y)$ pour cette
énumération. Pourquoi $g(x,y)$ est-elle calculable ?
Intuitivement, la réponse est claire : pour calculer $g(x,y)$,
on décode d'abord~$x$ et on vérifie s'il s'agit d'une
notation de fonction unaire. Si c'est le cas, on calcule
la valeur de cette fonction en~$y$; sinon, on renvoie~$0$.

\begin{tagblock}{TMs}
\begin{digress}
Vous êtes peut-être déjà convaincu qu'avec un certain travail,
on peut écrire un programme, en Java ou en C++ par exemple,
qui effectue ce calcul. Nous pouvons alors invoquer la thèse
de Church--Turing : tout ce qui est calculable au sens intuitif
peut être calculé par une machine de Turing.

Une démonstration plus directe consisterait à décrire explicitement
une machine de Turing qui calcule~$g$. Elle éviterait tout recours
à la thèse de Church--Turing ou à l'intuition. Nous disposerons
bientôt des outils nécessaires pour montrer que $g$ est calculable
au moyen d'un modèle de calcul que l'on peut \emph{simuler} sur
une machine de Turing : celui des fonctions récursives.
\end{digress}
\end{tagblock}

\end{document}
